%%%%%%%%%%%%%%%%%%%%%%%%%%%%%%%%%%%%%%%%%%%%%%%%%%%%%%%
%                   Źródło: Candela_abstrakt.tex           %
%                  Data: 25 Marca 2022                %
%                                                     %
%     Plik do generowania abstraktu na Polską         %
%     Konferencję Optyczną 2022 w LaTeX               %
%                                                     %
%                                                     %
%       (c) 2021 Optica, 2023 Fundacja Candela         %
%%%%%%%%%%%%%%%%%%%%%%%%%%%%%%%%%%%%%%%%%%%%%%%%%%%%%%%

\documentclass[letterpaper, 10pt]{article}
%% dostępne kartki: A4 lub letterpaper

\usepackage{candela3} %% Proszę nie zmieniać
\newgeometry{tmargin=1.75cm, bmargin=1.75cm, lmargin=1.75cm, rmargin=1.75cm}

%% tu proszę ładować ew. dodatkowe biblioteki. Proszę nie usuwać
%% tych załadowanych domyślnie.

\usepackage[colorlinks=true,bookmarks=false,citecolor=black,urlcolor=black]{hyperref} %pdflatex



\begin{document}

\title{Zgłoszenie na Polską Konferencję Optyczną 2023}  %% Proszę tutaj wpisać tytuł zgłoszenia

\author{Pierwszy Autor$^1$, Drugi Autor$^2$}

\address{
$^1$Afiliacja Autora z adresem\\
$^2$Afiliacja Autora z adresem}

\email{e-mail Autora}
\copyrightyear{2023}


\section{Treść zgłoszenia}

\subsection{Elementy wymagane}
Wszystkie przesłane pliki PDF zawierające zgłoszenie na Polską Konferencję Optyczną 2023 muszą spełniać poniższe wymagania:

\begin{enumerate}
\item zgłoszenie musi mieć tytuł
\item zgłoszenie musi zawierać kompletną listę jego autorów wraz z ich afiliacjami
\item zgłoszenie musi stanowić całość (część zgłoszenia nie może znajdować się poza przesłanym plikiem PDF)
\item zgłoszenie nie może zawierać elementów, na których wykorzystywanie Autorzy nie posiadają zgody (dotyczy to np. danych, grafik, logotypów itp.)
\item \textbf{maksymalna długość zgłoszenia to jedna strona}.
\end{enumerate}

\subsection{Styl typograficzny}
Marginesy, rozmiary czcionek, styl i umiejscowienie tytułu, imion autorów, adresów, treści, bibliografii, podpisów itp. są narzucone przez plik \texttt{candela3.sty} i nie należy ich zmieniać - to zapewni spójny wygląd wszystkich przyjętych zgłoszeń. Dopuszczalne jest dodatkowe użycie biblioteki \texttt{ae.sty}.

\subsection{Imiona i nazwiska Autorów oraz ich afiliacje}
Nazwiska autorów należy wpisać w całości, a imiona w postaci inicjałów (np. Jan Adam Kowalski $\rightarrow$ J. A. Kowalski). Afiliacje należy wpisać z pełnymi adresami i bez skrótów, przynajmniej na poziomie Wydziałów / Instytutów itp. (np. Wydział Fizyki Uniwersytetu Warszawskiego, ul. Pasteura 5, 02-093 Warszawa, Polska).

\subsection{Notacja}
\subsubsection{Oznaczenia}
Użyte oznaczenia powinny być czytelne, jasne, zwięzłe i najlepiej zgodne z tymi stosowanymi powszechnie. Akronimy są dopuszczalne, o ile zdefiniuje się ich znaczenie przy pierwszym użyciu.

\subsubsection{Notacje matematyczne}
Do równań należy stosować standardowe komendy \LaTeX{} lub AMS\LaTeX{}. Poniżej znajduje się przykład z Krishnan \textit{et al.} \cite{krishnan00}.

\begin{eqnarray}
\bar\varepsilon &=& \frac{\int_0^\infty\varepsilon
\exp(-\beta\varepsilon)\,{\rm d}\varepsilon}{\int_0^\infty
\exp(-\beta\varepsilon)\,{\rm d}\varepsilon}\nonumber\\
&=& -\frac{{\rm d}}{{\rm d}\beta}\log\Biggl[\int_0^\infty\exp
(-\beta\varepsilon)\,{\rm d}\varepsilon\Biggr]=\frac1\beta=kT.
\end{eqnarray}

\section{Tabele i rysunki}
Rysunki, wykresy, ilustracje i tabele należy umieścić bezpośrednio w przesłanym pliku PDF. Ich rozmiar powinien być adekwatny do wartości informacji zawartej na danym rysunku. Maksymalna liczba rysunków to 3, przy czym każdy z nich powinien być umiejscowiony możliwie blisko tej części tekstu, w której jest omawiany. Szerokość rysunków nie może przekraczać szerokości tekstu.

\begin{figure}[htbp]
  \centering
  \includegraphics[width=5cm]{OT10000F1.eps}
\caption{Przykładowy rysunek wraz z preferowanym stylem podpisu.}
\end{figure}


\begin{table}[htb]
 \centering \caption{Przykładowa tabela.}
\begin{tabular}{ccc}
    \hline
    Jeden & Dwa & Trzy \\
    \hline
    One & Two & Three \\
    Un & Deux & Trois \\
    Jeden & Dv\v{e} & T\v{r}i \\
    \hline
   \end{tabular}
    \end{table}

\section{Literatura}
Odniesienia literaturowe proszę cytować za pomocą komendy ``cite``, która wykorzystuje preferowany format cytowania w nawiasach kwadratowych \cite{krishnan00,vantrigt97,masters93,shoop97,kalman76,craig96,steup96}. Styl \texttt{candela3.sty} korzysta z cytowań jak w \texttt{cite.sty}. Więcej informacji na temat sposobu cytowania znaleźć można na stronie \href{http://www.crossref.org/}{http://www.crossref.org/}.

\begin{thebibliography}{99} %% można albo wykorzystać bibliotekę BibTeX albo ręcznie dodać odniesienia zgodnie z przykładami poniżej:

\bibitem{krishnan00} E. Krishnan, A. M. Shan, T. Rishi, L. A. Ajith, C. V.
Radhakrishnan, \textit{On-line Tutorial on \LaTeX{}},
``Mathematics'' (Indian \TeX{} Users Group, 2000), \\
\url{http://www.tug.org/tutorials/tugindia/chap11-scr.pdf}.

\bibitem{vantrigt97} C. van Trigt, ``Visual system-response functions and estimating reflectance,''
J. Opt. Soc. Am. A \textbf{14}, 741--755 (1997).

\bibitem{masters93} T. Masters, \emph{Practical Neural Network Recipes in C++} (Academic, 1993).

\bibitem{shoop97} B. L. Shoop, A. H. Sayles, and D. M. Litynski, ``New devices for optoelectronics: smart pixels,''
in \emph{Handbook of Fiber Optic Data Communications},
C. DeCusatis, D. Clement, E. Maass, and R. Lasky, eds. (Academic, 1997), pp. 705--758.

\bibitem{kalman76} R. E. Kalman,``Algebraic aspects of the generalized inverse of a rectangular matrix,'' in
\emph{Proceedings of Advanced Seminar on Generalized Inverse and Applications}, M. Z. Nashed, ed. (Academic, 1976), pp. 111--124.

\bibitem{craig96} R. Craig and B. Gignac, ``High-power 980-nm pump lasers,''
in \emph{Optical Fiber Communication Conference}, Vol. 2 of 1996 OSA Technical Digest Series (Optical Society of America, 1996), paper ThG1.

\bibitem{steup96} D. Steup and J. Weinzierl, ``Resonant THz-meshes,''
presented at the Fourth International Workshop on THz Electronics, Erlangen-Tennenlohe, Germany, 5--6 Sept. 1996.

\end{thebibliography}

\end{document}
